\documentclass[12pt]{article}
\usepackage[utf8]{inputenc}
\usepackage[T1]{fontenc}
\usepackage{hyperref}
\title{Zapotec Language Resources for K-12 Curriculum Development}
\author{Compiled by Claude Sonnet 5.5 (claude-sonnet-5-5) for Andr\'as Kornai}
\date{}
\begin{document}
\maketitle

\section{Introduction}
The code zap is the ISO 639-3 macrolanguage Zapotec, which covers many mutually distinct languages of Oaxaca, Mexico. The resources below concern particular varieties (mainly Valley Zapotec and a Sierra Norte variety), not the whole macrolanguage, and only some of the member languages are covered here.

\section{Primary Educational Resources}
\begin{itemize}
\item \emph{Cali Chiu: A Course in Valley Zapotec} (3rd ed.), an open educational resource by Pamela Munro, Brook Danielle Lillehaugen and Felipe H. Lopez, with lessons, videos, cultural material and a dictionary: \url{https://scholarship.haverford.edu/books/67}
\item Study on communicative competence and motivation among learners of Zapotec in Teotitl\'an del Valle: \url{https://scholarspace.manoa.hawaii.edu/handle/10125/25380?mode=full}
\end{itemize}

\section{Linguistic and Reference Materials}
The Cali Chiu course above includes a reference section and dictionary. No other grammars were reviewed.

\section{Religious and Community Resources}
UC Santa Cruz's Zapotec Language Project, working with a Sierra Norte variety since 2016, and its partner Senderos (Nido de Lenguas): \url{https://zapotec.ucsc.edu/history.php}

\section{Digital and Audio Resources}
The Cali Chiu course embeds videos. No separate audio archive was found for the macrolanguage.

\section{Curriculum Development Context}
Teaching Zapotec is mostly done through university programmes and community classes, and one source notes a lack of second-language teaching materials and instructor training. Any curriculum must be built for one specific variety, and the individual-language codes (for example zam and zao) should be searched separately.

\end{document}